\setcounter{section}{1}
\setcounter{subsection}{1}
\setcounter{theorem}{1}
\subsection{A Natural Complete Problem via Chevalley-Warning Theorem}
\label{sec:intro-complete}

There have been several works focusing on completeness results for the class $\PPA$ (i.e. $\PPA_2$). Initial works showed the
$\PPA$-completeness of (non-natural) total search problems corresponding to topological fixed point theorems
\cite{r24,r1,r19}. Closer to our paper,
Belovs et al. \cite{r8} show the $\PPA$-completeness of computational analogs of Combinatorial Nullstellensatz and the Chevalley--Warning Theorem, but which explicitly
involve a circuit as part of the input. More recently, breakthrough results
showed $\PPA$-completeness of problems without a circuit or a Turing Machine in
the input such as $\textsc{Consensus-Halving}$, $\textsc{Necklace-Splitting}$
and $\textsc{Ham-Sandwich}$
\cite{r20,r21} resolving an open
problem since the definition of $\PPA$ in \cite{r33}.

  Our main contribution is to provide a natural complete problem for $\PPA_p$,
for every prime $p$; thereby also yielding a new complete problem for $\PPA$.
Our complete problem is an extension of the problem $\eChevalleyp$, defined by
Papadimitriou \cite{r33}, which is a search problem
associated to the celebrated Chevalley-Warning Theorem. We first
present an abstract way to understand the proof of the Chevalley-Warning Theorem
that motivates the definition of our natural complete problem for
$\PPA_p$.

\subsubsection{Max-Degree Monic Monomials and Proof of Chevalley-Warning Theorem} \label{sec:maximum-degree}

  In 1935, Claude Chevalley \cite{r15} resolved a
hypothesis stated by Emil Artin, that all finite fields are quasi-algebraically
closed. Later, Ewald Warning~\cite{r37}
proved a slight generalization of Chevalley's theorem. This
generalized statement is usually referred to as the Chevalley-Warning Theorem
(\CWT, for short). Despite its initial algebraic motivation, \CWT has found
profound applications in combinatorics and number theory as we discuss in
Section1.4 (and Section6).

We now explain the statement of the Chevalley-Warning Theorem, starting with some notations. For any field $\bbF$ and any polynomial $f$ in a polynomial ring $\bbF[x_1, \ldots, x_n]$ we use $\deg(f)$ to represent the degree of $f$. We use $\vecx$ to succinctly denote the set of all variables $(x_1, \ldots, x_n)$ (the number of variables will always be $n$) and $\vecf$ to succinctly denote a system of polynomials $\vecf = (f_1, \ldots, f_m) \in \bbF[\vecx]^m$. We will often abuse notations to use $\bx$ to also denote assignments over $\bbF_p^n$. For instance, let $\calV_{\vecf} \coloneqq \set{\bx \in \bbF_p^n : f_i(\bx) = 0 \text{ for all } i \in [m]}$ be the set of all common roots of $\vecf$.

\theoremstyle{plain}
\newtheorem*{cwthm}{Chevalley-Warning Theorem}
\begin{cwthm}[\cite{r15,r37}]
For any prime\footnote{While most of the results in this section
generalize to prime powers, we only consider prime fields for simplicity.} $p$ and polynomial system
$\vecf \in \bbF_p[\vecx]^m$ satisfying\vspace{-3.5mm}
\begin{equation}
\tag{CW Condition}
\sum_{i = 1}^m \deg(f_i) < n,\vspace*{-3mm}
\label{eq:CWcondition}
\end{equation}
it holds that $|\calV_{\vecf}| \equiv 0 \mod{p}$.
\end{cwthm}

\noindent Given a polynomial system $\vecf \in \bbF_p[\vecx]^m$,
the key idea in the proof of the Chevalley-Warning Theorem is the polynomial
\[ \CW_{\vecf}(\bx) \coloneqq \prod_{i = 1}^m \inparen{1 - f_i(\vecx)^{p - 1}} \mod{\set{x_i^p - x_i}_i}\,. \] Observe that $\CW_{\vecf}(\vecx) = 1$ if $\vecx \in \calV_{\vecf}$ and is $0$ otherwise. Thus, $|\calV_{\vecf}| \equiv \sum_{\vecx \in \bbF_p^n} \CW_{\vecf}(\bx) \mod{p}$.
The following definition informally describes a special type of monomial of
$\CW_{\vecf}$ that is of particular interest in the proof. For the precise definition, we
refer to \autoref{sec:chevalley}.

\begin{definition}[\textsc{Max-Degree Monic Monomials} (Informal)]
	\label{def:max-degreeMonomialInformal}
    Let $\vecf \in \bbF_p[\vecx]^m$.
    			A \textit{monic monomial} of $\CW_{\vecf}$ refers to a monic monomial obtained when symbolically expanding $\CW_{\vecf}$ as a sum of monic monomials. A monic monomial is said to be of \textit{max-degree} if it is
	$\prod_{j = 1}^n x_j^{p - 1}$.

\end{definition}

   In the above definition, it is important to consider the
\emph{symbolic expansion} of $\CW_{\vecf}$ and ignore any cancellation of coefficients that might occur. Observe that, although the
expansion of $\CW_{\vecf}$ is exponentially large in the description size of
$\vecf$, each monic monomial of $\CW_{\vecf}$ can be succinctly described as a combination of monic monomials of the polynomials $f_1, \ldots, f_m$. We formally
discuss this in \autoref{sec:chevalley}.

Using the definition of max-degree monic monomials, we state the main technical
lemma underlying the proof of \CWT (with proof in \autoref{sec:chevalley}).
\theoremstyle{plain}
\newtheorem*{cwthmprinciple}{Chevalley--Warning Lemma}
\begin{cwthmprinciple}
For any prime $p$ and $\vecf \in \bbF_p[\vecx]^m$,
\begin{equation}
	\tag{CW Lemma}
	|\calV_{\vecf}| ~\equiv~ (-1)^n \cdot  |\set{\mathrm{max\text{-}degree~monic~monomials~of~} \CW_{\vecf}}| \mod{p}
	\label{eq:mainPrinciple}
\end{equation}
\end{cwthmprinciple}

\noindent The Chevalley-Warning Theorem now follows by observing that if $\sum_{i = 1}^m \deg(f_i) < n$ then the number of max-degree monic monomials of $\CW_{\vecf}$ is zero. Hence, we get that $|\calV_{\vecf}| \equiv 0 \mod{p}$.

\subsubsection{Proofs of Cancellation}
\label{sec:proofsOfCancellation}

  From the proof sketch of CWT in the previous section, a slight generalization of CWT follows. In particular, $|\calV_{\vecf}| \equiv 0 \mod{p}$ if and only if
\begin{equation}
	\tag{Extended CW Condition}
	\inabs{\set{\mathrm{max\text{-}degree~monic~monomials~of~} \CW_{\vecf}}} ~\equiv~ 0 \mod{p},
		\label{eq:generalizedCondition}
\end{equation}
Thus, any condition on $\vecf$ that implies the \eqref{eq:generalizedCondition} can
 replace \eqref{eq:CWcondition} in the Chevalley-Warning
Theorem. Note that the \eqref{eq:generalizedCondition} is equivalent to all the max-degree monic monomials in $\CW_{\vecf}$ cancelling out.
Thus, we call any such condition on $\vecf$ that implies \eqref{eq:generalizedCondition} to be a
``\textit{proof of cancellation}'' for the system $\vecf$.

We can now reinterpret the result of Belovs et
al. \cite{r8} in this framework of ``proof of cancellation'' conditions. In particular, \cite{r8} considers the
case $p = 2$ and defines the problem $\textsc{PPA-Circuit-Chevalley}$, in which a ``proof of
cancellation'' is given in a specific form of circuits. These circuits describe the
system $(f_1, \ldots, f_m)$ in the $\textsc{PPA-Circuit-Chevalley}$ problem. It is then shown that $\textsc{PPA-Circuit-Chevalley}$ is $\PPA_2$-complete.


\subsubsection{Computational Problems Based on Chevalley-Warning Theorem}
\label{sec:computationalChevalleyWarning}

  Every ``proof of cancellation'' that is \textit{syntactically refutable} can be used to define
  a total search problem that lies in $\PPA_p$. By {\em syntactically refutable} we mean that whenever the
  ``proof of cancellation'' is false, there exists a small witness that certifies so. In this
  section, we define three computational problems with their corresponding ``proof of
  cancellation'': (1) the $\eChevalleyp$ problem defined by \cite{r33}, (2) the
  $\lChevalleyp$ problem that is a generalization of $\eChevalleyp$, and (3) the problem
  $\cChevalleyp$ that we show to be $\PPA_p$-complete. All these problems are
  defined for every prime modulus $p$ and are natural in the sense that they do not explicitly
  involve a circuit or a Turing Machine in their input. In particular, the polynomial systems in the input are {\em explicit} in that they are given as a sum of monic monomials. \bigskip

\paragraph{Chevalley.}
  This is the direct computational analog of the Chevalley-Warning Theorem and was
defined by Papadimitriou \cite{r33} as the following total
search problem:
\medskip

\noindent $\underbar{\eChevalley}_{\pnum}$

\noindent Given an explicit polynomial system $\vecf \in \bbF_p[\vecx]^m$, and an
$\vecx^{\star} \in \calV_{\vecf}$, output one of the following:\vspace*{-0.8mm}
\begin{itemize}
	\item[$\triangleright$] [{\em Refuting witness}] \eqref{eq:CWcondition} is not satisfied.
	\item[$\triangleright$] $\vecx \in \calV_{\vecf} \smallsetminus \set{\vecx^{\star}}$.
\end{itemize}

\noindent We will particularly consider a special case where all the $f_i$'s have zero constant term
({\em zecote}, for short). In this case, $\vecx^{\star} = \mathbf{0} \in \calV_{\vecf}$, so there is no need to explicitly include $\vecx^*$ in the input.

\paragraph{General Chevalley.} As mentioned already, we can define a search problem corresponding to any syntactically refutable condition that implies the
\eqref{eq:generalizedCondition}. One such condition is to directly assert that\vspace{-2.5mm}
\begin{equation}
	\tag{General CW Condition}
	\{\mathrm{max\text{-}degree~monic~monomials~of~} \CW_{\vecf}\} ~=~ \emptyset. \label{eq:generalCWCondition}
\end{equation}
In particular, note that \eqref{eq:CWcondition} implies this condition. Moreover, this condition is syntactically refutable by a max-degree monic monomial, which is efficiently representable as a combination of at most $m(p-1)$ monomials of the $f_i$'s. Thus, we can define the following total search problem generalizing $\eChevalleyp$.\medskip

\noindent $\underbar{\lChevalley}_{\pnum}$

\noindent Given an explicit polynomial system $\vecf \in \bbF_p[\vecx]^m$ and an
$\vecx^{\star} \in \calV_{\vecf}$, output one of the following:\vspace{-0.5mm}
\begin{itemize}
	\item[$\triangleright$] [{\em Refuting Witness}] A max-degree monic monomial of $\CW_{\vecf}$.
	\item[$\triangleright$] $\vecx \in \calV_{\vecf} \smallsetminus \set{\vecx^{\star}}$.
\end{itemize}
\medskip

\noindent While $\lChevalleyp$ generalizes $\eChevalleyp$, it does not
capture the full generality of \eqref{eq:generalizedCondition}. However \eqref{eq:generalizedCondition} is not syntactically refutable (in fact, it is $\mathsf{Mod}_p\mathsf{P}$--complete to decide\footnote{Circuit-SAT can be encoded as satisfiability of a polynomial system $\vecf \in \bbF_p[\vecx]^m$ by including a polynomial for each gate along with $\set{x_i^2 - x_i = 0}$ to ensure Booleanity. Thus, number of satisfiable assignments to the Circuit-SAT is $\equiv |\calV_{\vecf}| \mod{p}$, which is $0 \mod{p}$ iff the final coefficient of the max-degree monomial is $0$.} if the final coefficient of the max-degree monomial is $0$).

A natural question then is whether $\lChevalley_p$, or even $\Chevalley_p$, could already be $\PPA_p$--complete. We believe this to be unlikely because \eqref{eq:generalCWCondition} seems to fail in capturing other simple conditions that are syntactically refutable and yet imply \eqref{eq:generalizedCondition}. Namely, consider a permutation $\sigma \in S_n$ of the variables $x_1, \ldots, x_n$ of order $p$ (i.e. $\sigma^p$ is the identity permutation). Suppose that for every $\vecx \in \overline{\calV_{\vecf}}$, it holds that $\sigma(\vecx) \in \overline{\calV_{\vecf}} \smallsetminus \set{\vecx}$; in other words $\vecx, \sigma(\vecx), \sigma^2(\vecx), \ldots, \sigma^{p-1}(\vecx)$ are all distinct and in $\overline{\calV_{\vecf}}$ (where, $\sigma(\vecx)$ denotes the assignment obtained by permutating the variables of the assignment $\vecx$ according to $\sigma$); observe that this condition is syntactically refutable. This implies that the elements of $\overline{\calV_{\vecf}}$ can be partitioned into groups of size $p$ (given by the orbits of the action $\sigma$) and hence $|\overline{\calV_{\vecf}}| \equiv 0 \mod{p}$. Thus, such a $\sigma$ provides a syntactically refutable proof that $|\calV_{\vecf}| \equiv 0 \mod{p}$ and hence that \eqref{eq:generalizedCondition} hold.

Hence, we further generalize 
$\lChevalleyp$ into a problem that incorporates this additional ``proof of cancellation'' in the form of a permutation $\sigma \in S_n$.

\paragraph{Chevalley with Symmetry.}
We consider a union of two polynomial systems $\vecg  \in \bbF_p[\vecx]^{m_g}$ and $\vech \in \bbF_p[\vecx]^{m_h}$. Even if both $\vecg$ and
$\vech$ satisfy \eqref{eq:CWcondition}, the combined system $\vecf :=
(g_1, \ldots, g_{m_g}, h_1, \ldots, h_{m_h})$ might not satisfy \eqref{eq:CWcondition}
and it might even be the case that $|\calV_{\vecf}|$ is not a multiple of $p$. Thus, we need to bring in some additional conditions.

We start by observing that since $|\calV_{\vecf}| + |\overline{\calV_{\vecf}}| = p^n$, it holds that $|\calV_{\vecf}| \equiv 0 \mod{p}$ if and only if $|\overline{\calV_{\vecf}}| \equiv 0 \mod{p}$. Also note that, $|\overline{\calV_{\vecf}}| = |\overline{\calV_{\vecg}}| + |(\calV_{\vecg} \cap \overline{\calV_{\vech}})|$.


If $\vecg$ satisfies the \eqref{eq:generalCWCondition} then we have that $|\calV_{\vecg}| \equiv |\overline{\calV_{\vecg}}| \equiv 0 \mod{p}$. A simple way to enforce that $|\calV_{\vecg} \cap \overline{\calV_{\vech}}| \equiv 0 \mod{p}$ is to enforce a ``symmetry'', namely that its elements can be grouped into groups of size $p$ each. We impose this grouping with a permutation $\sigma \in S_n$ of the variables $x_1, \ldots, x_n$ of order $p$ such that for any $\vecx \in \calV_{\vecg} \cap \overline{\calV_{\vech}}$, it holds that $\sigma(\vecx) \in (\calV_{\vecg} \cap \overline{\calV_{\vech}}) \smallsetminus \set{\vecx}$; or in other words that $\vecx, \sigma(\vecx), \sigma^2(\vecx), \ldots, \sigma^{p-1}(\vecx)$ are all distinct and contained in $\calV_{\vecg} \cap \overline{\calV_{\vech}}$.

We now define the
following natural total search problem.

\bigskip

\noindent $\underbar{\cChevalley}_{\pnum}$

\noindent Given two explicit polynomial systems $\vecg \in \bbF_p[\vecx]^{m_g}$ and $\vech \in \bbF_p[\vecx]^{m_h}$, and an $\vecx^{\star} \in \calV_{\vecf}$ (where $\vecf \coloneqq (\vecg, \vech)$) and a
permutation $\sigma \in S_n$ of order $p$, output one of the following:\vspace{-0.5mm}
\begin{itemize}
	\item[$\triangleright$] [{\em Refuting Witness} -- 1] A max-degree monic monomial of $\CW_{\vecg}$.
  \item[$\triangleright$] [{\em Refuting Witness} -- 2] $\vecx \in \calV_{\vecg} \cap \overline{\calV_{\vech}}$ such that $\sigma(\vecx) \notin (\calV_{\vecg} \cap \overline{\calV_{\vech}}) \smallsetminus \set{\vecx}$.
	\item[$\triangleright$] $\vecx \in \calV_{\vecf} \smallsetminus \set{\vecx^{\star}}$.
\end{itemize}
\medskip

\noindent The above problem is natural, because the input consists of a system of polynomial in an
explicit form, i.e. as a sum of monic monomials, together with a permutation in $S_n$ given say in
one-line notation. Also, observe that when $\vech$ is empty, the above problem coincides with $\lChevalley_p$ (since $\overline{\calV_{\vech}} = \emptyset$ when $\vech$ is empty). Our main result is the following (proved in \autoref{sec:chevalley}).
